% \lstdefinelanguage{Domino}{
%     % morecomment=[s]{/*}{*/},
%     moredelim=[is][\color{gray}]{@@}{@@}
% }
% \begin{lstlisting}[
%     language=Domino,
%     caption={Example Domino theorem},
%     label={lst:domino-theorem-example},
%     commentstyle=\small\rmfamily\itshape,
%     basicstyle=\small\ttfamily,
%     columns=fixed]
%     theorem Theorem1 {
%     @@/*an integer const parameter passed @@
%     @@ to game instances*/@@
%     const n: Integer;
%     @@/*a const function that is passed to @@
%     @@ game instances*/@@
%     const fn1: fn Bits(n), Bits(n), -> Bits(m);
%     @@/*Game instances used in the theorem*/@@
%     instance LeftGame = Game_R_Pkg1{
%         params { n: n, fn1: fn1 }
%     }
%     instance RightGame = Game2{
%         params { n:n, fn1: fn1 }
%     }
%     instance BaseLGame = Game3{
%         params { n:n, fn1: fn1 }
%     }
%     instance BaseRGame = Game4{
%         params { n:n, fn1: fn1 }
%     }
%     @@/*... other games...*/@@
%     assumptions {BaseGame: BaseLGame ~ BaseRGame}
%     gamehops {
%         reduction LeftGame RightGame {
%             assumption BaseGame
%             map {
%                 @@/*map packages of LeftGame and BaseLGame*/@@
%             }
%             map {
%                 @@/*map packages of RightGame and BaseRGame*/@@
%             }
%         }
%         equivalence GameN GameM {
%             invariant: [
%                 ./theorem/invariants/invariant.smt2
%             ]
%             ORACLE: {
%                 lemmas {
%                     invariant: [no-abort]
%                     same-output: [no-abort]
%                     equal-aborts: []
%                     @@/*other lemmata*/@@
%                 }
%             }
%         }
%         @@...@@
%     }}
% \end{lstlisting}
\begin{longlisting}
\begin{minted}{ssp-lang}
theorem Theorem1 {
    /*an integer const parameter passed 
     to game instances*/
    const n: Integer;
    /*a const function that is passed to 
     game instances*/
    const fn1: fn Bits(n), Bits(n), -> Bits(m);
    /*Game instances used in the theorem*/
    instance LeftGame = Game_R_Pkg1{
        params { n: n, fn1: fn1 }
    }
    instance RightGame = Game2{
        params { n:n, fn1: fn1 }
    }
    instance BaseLGame = Game3{
        params { n:n, fn1: fn1 }
    }
    instance BaseRGame = Game4{
        params { n:n, fn1: fn1 }
    }
    /*... other games...*/
    assumptions {BaseGame: BaseLGame ~ BaseRGame}
    gamehops {
    reduction LeftGame RightGame {
        assumption BaseGame
        map {
        /*map packages of LeftGame and BaseLGame*/
        }
        map {
        /*map packages of RightGame and BaseRGame*/
        }
    }
    equivalence GameN GameM {
        invariant: [
            ./theorem/invariants/invariant.smt2
        ]
        ORACLE: {
            lemmas {
                invariant: [no-abort]
                same-output: [no-abort]
                equal-aborts: []
                /*other lemmata*/
            }
        }
    }
    /*...*/
}}
\end{minted}
\caption{Example Domino theorem}
\label{lst:domino-theorem-example}
\end{longlisting}